\documentclass[11pt,a4paper]{article}
% preamble.tex — shared setup for human-verification documents
% Compile each document with:  tectonic <File>.tex

\usepackage[T1]{fontenc}
\usepackage[utf8]{inputenc}
\usepackage{lmodern}
\usepackage{amssymb,amsmath,amsthm}
\usepackage{mathtools}
\usepackage[dvipsnames]{xcolor}
\usepackage{listings}
\usepackage{hyperref}
\usepackage[margin=2.5cm]{geometry}
\usepackage{parskip}

% ── Colours ──────────────────────────────────────────────────────────
\definecolor{leanBg}{HTML}{F7F7F0}
\definecolor{haskellBg}{HTML}{F0F4F7}
\definecolor{keyword}{HTML}{7B3294}
\definecolor{comment}{HTML}{5A8C32}
\definecolor{noteFg}{HTML}{B35900}

% ── Lean 4 listings language ─────────────────────────────────────────
% NOTE: Lean listing blocks use ASCII approximations of unicode symbols
% (e.g. -> for arrow, o. for compose).  The literate key substitutes
% these with proper math-mode renderings in the PDF.
\lstdefinelanguage{Lean4}{
  morekeywords={def,theorem,structure,where,namespace,end,import,fun,
                let,have,show,by,exact,intro,cases,rfl,subst,obtain,
                funext,simp,rw,apply,variable,Type,Prop,instance,
                section,open,class,extends,deriving,sorry,noncomputable},
  sensitive=true,
  morecomment=[l]{--},
  morecomment=[n]{/-}{-/},
  morestring=[b]",
}

% ── Lean listing style ───────────────────────────────────────────────
\lstdefinestyle{lean}{
  language=Lean4,
  backgroundcolor=\color{leanBg},
  basicstyle=\small\ttfamily,
  keywordstyle=\color{keyword}\bfseries,
  commentstyle=\color{comment}\itshape,
  breaklines=true,
  frame=single,
  framerule=0.4pt,
  xleftmargin=4pt,
  xrightmargin=4pt,
  aboveskip=6pt,
  belowskip=4pt,
  literate=
    {->}{{$\to$}}2
    {=>}{{$\Rightarrow$}}2
    {<|}{{$\langle$}}1
    {|>}{{$\rangle$}}1
    {==>}{{$\Longrightarrow$}}3
    { o. }{{$\;\circ\;$}}3
    {forall}{{$\forall$}}1
    {:=}{{$\,:=\,$}}2
    {!=}{{$\neq$}}2
    {<=}{{$\leq$}}2
    {>=}{{$\geq$}}2,
}

% ── Haskell listing style ───────────────────────────────────────────
\lstdefinestyle{haskell}{
  language=Haskell,
  backgroundcolor=\color{haskellBg},
  basicstyle=\small\ttfamily,
  keywordstyle=\color{keyword}\bfseries,
  commentstyle=\color{comment}\itshape,
  breaklines=true,
  frame=single,
  framerule=0.4pt,
  xleftmargin=4pt,
  xrightmargin=4pt,
  aboveskip=6pt,
  belowskip=4pt,
  literate=
    {->}{{$\to$}}2
    {=>}{{$\Rightarrow$}}2,
}

% ── Theorem environments ────────────────────────────────────────────
\theoremstyle{definition}
\newtheorem{defn}{Definition}[section]
\newtheorem{thm}[defn]{Theorem}
\newtheorem{prop}[defn]{Proposition}

\theoremstyle{remark}
\newtheorem*{fidelitynote}{Fidelity note}

% ── Convenience macros ───────────────────────────────────────────────
\newcommand{\lean}[1]{\texttt{#1}}      % inline Lean identifier
\newcommand{\hask}[1]{\texttt{#1}}      % inline Haskell identifier
\newcommand{\Ext}{\mathrm{Ext}}
\newcommand{\Pos}{\mathrm{Pos}}
\newcommand{\Shape}{\mathrm{Shape}}


\title{Human Verification: \texttt{Containers.Tracking}\\[4pt]
  \large Compositional Tracking on Shipping Networks}
\author{Machine-generated from \texttt{lean/Tracking.lean} for review by Neil Ghani}
\date{July 2026}

\begin{document}
\maketitle

\begin{abstract}
This document presents every definition and theorem \emph{statement} from the
Lean~4 module \lean{Containers.Tracking} in three parallel forms: the original
Lean, a readable Haskell rendering, and precise mathematical English.  Proofs
are omitted --- they are machine-checked by Lean.  The purpose is to let a
human verifier confirm that the formalisation captures the intended
mathematics.

The file verifies Construction~\texttt{con:tracking} (\S4.3 of the paper): for
any monoid $M$, a \textbf{compositional tracking} of $M$-data on a shipping
network is in bijection with edge-labellings $G.E \to M$.  The file also
provides three concrete monoid instances and a cold-chain worked example.
\end{abstract}

\tableofcontents
\bigskip

%===================================================================
\section{Overview}
%===================================================================

A tracking assigns a monoid value $d(v, \rho)$ to every route $\rho$ out of
node $v$, compatibly with the directed-container structure:
\begin{itemize}
  \item \textbf{T-root}: $d(v, \mathrm{root}(v)) = 1$ (empty route gets unit);
  \item \textbf{T-shift}: $d(v, \mathrm{shift}(v, \rho, \sigma)) = d(v, \rho)
    \cdot d(\mathrm{sub}(v, \rho), \sigma)$ (concatenation becomes product).
\end{itemize}
The main theorem says this is exactly an edge-labelling: fold the label along the
path.


%===================================================================
\section{Bundled monoid}
%===================================================================

% --- Monoid' -------------------------------------------------------
\subsection{Structure: \lean{Monoid'}}

\begin{defn}[Bundled monoid]
A \textbf{bundled monoid} $M$ consists of a carrier type $M.\mathrm{M}$, a
binary operation $M.\mathrm{mul}$, a unit $M.\mathrm{one}$, and proofs of
associativity, left-unit, and right-unit.
\end{defn}

\paragraph{Lean.}
\begin{lstlisting}[style=lean]
structure Monoid' where
  M         : Type
  mul       : M -> M -> M
  one       : M
  mul_assoc : forall a b c, mul (mul a b) c = mul a (mul b c)
  one_mul   : forall a, mul one a = a
  mul_one   : forall a, mul a one = a
\end{lstlisting}

\paragraph{Haskell.}
\begin{lstlisting}[style=haskell]
-- Standard Haskell typeclass:
class Monoid m where
  mempty  :: m
  mappend :: m -> m -> m
  -- laws: mappend mempty x = x
  --       mappend x mempty = x
  --       mappend (mappend x y) z = mappend x (mappend y z)
\end{lstlisting}


%===================================================================
\section{The Tracking structure}
%===================================================================

% --- Tracking ------------------------------------------------------
\subsection{Structure: \lean{Tracking}}

\begin{defn}[Compositional tracking]
A \textbf{compositional tracking} of \lean{Mon}-data on the network
$\mathcal{N}(G)$ is a structure consisting of:
\begin{itemize}
  \item a function $d : (v : G.V) \to \mathcal{N}(G).\mathrm{Pos}(v) \to M$,
  \item proof of \textbf{T-root}: $d(v, \mathcal{N}(G).\mathrm{root}(v)) = 1$
    for all $v$, and
  \item proof of \textbf{T-shift}: for all $v$, $\rho$, $\sigma$,
    $d(v, \mathcal{N}(G).\mathrm{shift}(v, \rho, \sigma)) =
      d(v, \rho) \cdot d(\mathcal{N}(G).\mathrm{sub}(v, \rho), \sigma)$.
\end{itemize}
\end{defn}

\paragraph{Lean.}
\begin{lstlisting}[style=lean]
structure Tracking (G : Digraph) (Mon : Monoid') where
  d       : (v : G.V) -> (network G).Pos v -> Mon.M
  d_root  : forall v, d v ((network G).root v) = Mon.one
  d_shift : forall (v : G.V) (rho : (network G).Pos v)
              (sigma : (network G).Pos ((network G).sub v rho)),
      d v ((network G).shift v rho sigma) =
        Mon.mul (d v rho) (d ((network G).sub v rho) sigma)
\end{lstlisting}

\paragraph{Haskell.}
\begin{lstlisting}[style=haskell]
data Tracking g m = Tracking
  { trackVal   :: Node g -> Route g -> m
  -- laws (not enforced by the type):
  -- trackVal v (root v) = mempty
  -- trackVal v (shift v rho sigma) =
  --   mappend (trackVal v rho) (trackVal (sub v rho) sigma)
  }
\end{lstlisting}


%===================================================================
\section{Fold along a path}
%===================================================================

% --- foldPath ------------------------------------------------------
\subsection{Definition: \lean{foldPath}}

\begin{defn}[Path fold]
Given an edge-labelling $f : G.E \to M$, the \textbf{fold} of $f$ along a path
$\rho : \mathrm{Path}_G(v, w)$ is:
\[
  \mathrm{foldPath}(M, f, \mathrm{nil}(v)) = 1, \qquad
  \mathrm{foldPath}(M, f, \mathrm{cons}(e, \rho)) = f(e) \cdot \mathrm{foldPath}(M, f, \rho).
\]
\end{defn}

\paragraph{Lean.}
\begin{lstlisting}[style=lean]
def foldPath {G : Digraph} (Mon : Monoid') (f : G.E -> Mon.M) :
    {v w : G.V} -> Path G v w -> Mon.M
  | _, _, .nil _    => Mon.one
  | _, _, .cons e p => Mon.mul (f e) (foldPath Mon f p)
\end{lstlisting}

\paragraph{Haskell.}
\begin{lstlisting}[style=haskell]
foldPath :: Monoid m => (e -> m) -> [e] -> m
foldPath _ []     = mempty
foldPath f (e:es) = mappend (f e) (foldPath f es)
-- equivalently: foldMap f
\end{lstlisting}


% --- foldPath_append -----------------------------------------------
\subsection{Theorem: \lean{foldPath\_append}}

\begin{thm}[Fold distributes over concatenation]
For all edge-labellings $f$ and paths $\rho, \sigma$:
\[
  \mathrm{foldPath}(M, f, \rho \mathbin{+\!\!+} \sigma) =
  \mathrm{foldPath}(M, f, \rho) \cdot \mathrm{foldPath}(M, f, \sigma).
\]
\end{thm}

\paragraph{Lean.}
\begin{lstlisting}[style=lean]
theorem foldPath_append {G : Digraph} (Mon : Monoid') (f : G.E -> Mon.M) :
    forall {v w u : G.V} (p : Path G v w) (q : Path G w u),
    foldPath Mon f (p.append q) =
      Mon.mul (foldPath Mon f p) (foldPath Mon f q)
\end{lstlisting}

\paragraph{Haskell.}
\begin{lstlisting}[style=haskell]
-- foldPath f (p ++ q) = mappend (foldPath f p) (foldPath f q)
-- This is the standard homomorphism property of foldMap.
\end{lstlisting}


%===================================================================
\section{The bijection}
%===================================================================

% --- Tracking.ofEdge -----------------------------------------------
\subsection{Definition: \lean{Tracking.ofEdge}}

\begin{defn}[Edge-labelling $\to$ tracking]
Given an edge-labelling $f : G.E \to M$, the induced tracking has
$d(v, (w, \rho)) := \mathrm{foldPath}(M, f, \rho)$.  T-root holds because
$\mathrm{foldPath}$ of the empty path is $1$; T-shift holds by
\lean{foldPath\_append}.
\end{defn}

\paragraph{Lean.}
\begin{lstlisting}[style=lean]
def Tracking.ofEdge {G : Digraph} {Mon : Monoid'} (f : G.E -> Mon.M) :
    Tracking G Mon where
  d _ rho    := foldPath Mon f rho.2
  d_root _ := rfl
  d_shift _ rho sigma := foldPath_append Mon f rho.2 sigma.2
\end{lstlisting}

\paragraph{Haskell.}
\begin{lstlisting}[style=haskell]
ofEdge :: Monoid m => (e -> m) -> Tracking g m
ofEdge f = Tracking { trackVal = \_ route -> foldPath f (path route) }
\end{lstlisting}


% --- singleEdge / Tracking.toEdge ----------------------------------
\subsection{Definitions: \lean{singleEdge} and \lean{Tracking.toEdge}}

\begin{defn}[Single-edge position and restriction]
The \textbf{single-edge position} at the source of $e$ is the route
$(G.\mathrm{tgt}(e),\, \mathrm{cons}(e, \mathrm{nil}))
\in \mathcal{N}(G).\mathrm{Pos}(G.\mathrm{src}(e))$.

Given a tracking $t$, the \textbf{restriction to edges} is
$\mathrm{toEdge}(t)(e) := t.d(G.\mathrm{src}(e),\, \mathrm{singleEdge}(G, e))$.
\end{defn}

\paragraph{Lean.}
\begin{lstlisting}[style=lean]
def singleEdge (G : Digraph) (e : G.E) : (network G).Pos (G.src e) :=
  <|G.tgt e, Path.cons e (.nil (G.tgt e))|>

def Tracking.toEdge {G : Digraph} {Mon : Monoid'} (t : Tracking G Mon) :
    G.E -> Mon.M :=
  fun e => t.d (G.src e) (singleEdge G e)
\end{lstlisting}

\paragraph{Haskell.}
\begin{lstlisting}[style=haskell]
singleEdge :: Digraph v e -> e -> Route v e
singleEdge g e = (tgt g e, [e])

toEdge :: Tracking g m -> e -> m
toEdge t e = trackVal t (src g e) (singleEdge g e)
\end{lstlisting}


% --- round-trip theorems ------------------------------------------
\subsection{Theorems: \lean{Tracking.ofEdge\_toEdge} and \lean{Tracking.toEdge\_ofEdge}}

\begin{thm}[Bijection round trips]
The maps \lean{toEdge} and \lean{ofEdge} are mutually inverse:
\begin{enumerate}
  \item \textbf{Left inverse}: For every tracking $t$,
    $\mathrm{ofEdge}(\mathrm{toEdge}(t)) = t$.  (Uses T-root and T-shift
    induction along the path: \lean{foldPath\_eq\_tracking}.)
  \item \textbf{Right inverse}: For every edge-labelling $f$,
    $\mathrm{toEdge}(\mathrm{ofEdge}(f)) = f$.  (Immediate: fold of the
    single-edge path is $f(e) \cdot 1 = f(e)$.)
\end{enumerate}
\end{thm}

\paragraph{Lean.}
\begin{lstlisting}[style=lean]
theorem Tracking.ofEdge_toEdge {G : Digraph} {Mon : Monoid'}
    (t : Tracking G Mon) : Tracking.ofEdge (t.toEdge) = t

theorem Tracking.toEdge_ofEdge {G : Digraph} {Mon : Monoid'}
    (f : G.E -> Mon.M) : (Tracking.ofEdge f).toEdge = f
\end{lstlisting}

\paragraph{Haskell.}
\begin{lstlisting}[style=haskell]
-- ofEdge . toEdge = id    (for trackings)
-- toEdge . ofEdge = id    (for edge-labellings)
\end{lstlisting}


% --- trackingBijection ---------------------------------------------
\subsection{Definition: \lean{trackingBijection}}

\begin{defn}[Tracking bijection (Main Theorem)]
The pair $(\mathrm{toEdge}, \mathrm{ofEdge})$ constitutes a bijection
\[
  \mathrm{Tracking}(G, M) \;\cong\; (G.E \to M).
\]
Bundled as a \lean{Bijection} record with the two round-trip proofs.
\end{defn}

\paragraph{Lean.}
\begin{lstlisting}[style=lean]
def trackingBijection (G : Digraph) (Mon : Monoid') :
    Bijection (Tracking G Mon) (G.E -> Mon.M) where
  toFun     := Tracking.toEdge
  invFun    := Tracking.ofEdge
  left_inv  := Tracking.ofEdge_toEdge
  right_inv := Tracking.toEdge_ofEdge
\end{lstlisting}

\paragraph{Haskell.}
\begin{lstlisting}[style=haskell]
-- The bijection (toEdge, ofEdge) establishes:
-- Tracking G M  ~=  (Edge G -> M)
trackingBijection :: Iso (Tracking g m) (e -> m)
trackingBijection = Iso toEdge ofEdge
\end{lstlisting}


%===================================================================
\section{Monoid instances}
%===================================================================

% --- maxTempMonoid -------------------------------------------------
\subsection{Definition: \lean{maxTempMonoid}}

\begin{defn}[Max-temperature monoid]
The \textbf{max-temperature monoid} $(\{\mathrm{Cold}, \mathrm{Warm}\},
\max, \mathrm{Cold})$ is the idempotent monoid for peak-temperature tracking:
\lean{Warm} dominates \lean{Cold} (absorbs it), and \lean{Cold} is the unit.
A single warm leg along any route dominates all cold legs; no subsequent
refrigeration can undo it.
\end{defn}

\paragraph{Lean.}
\begin{lstlisting}[style=lean]
inductive MaxTemp | Cold | Warm
  deriving DecidableEq, Repr

def maxTempMonoid : Monoid' where
  M         := MaxTemp
  mul       := maxTempMul
  one       := .Cold
  one_mul   := fun x => by cases x <;> rfl
  mul_one   := fun x => by cases x <;> rfl
  mul_assoc := fun x y z => by cases x <;> cases y <;> cases z <;> rfl
\end{lstlisting}

\paragraph{Haskell.}
\begin{lstlisting}[style=haskell]
data MaxTemp = Cold | Warm deriving (Eq, Show)

instance Semigroup MaxTemp where
  Cold <> x = x
  Warm <> _ = Warm

instance Monoid MaxTemp where
  mempty = Cold
\end{lstlisting}


% --- coldChainMonoid -----------------------------------------------
\subsection{Definition: \lean{coldChainMonoid}}

\begin{defn}[Cold-chain monoid]
The \textbf{cold-chain monoid} $(\mathrm{Bool}, \mathbin{\&\&}, \mathrm{true})$:
$\mathrm{true}$ means the chain is intact; $\mathrm{false}$ means it is broken.
Once one leg breaks the chain, no subsequent leg can repair it.
\end{defn}

\paragraph{Lean.}
\begin{lstlisting}[style=lean]
def coldChainMonoid : Monoid' where
  M         := Bool
  mul       := (· && ·)
  one       := true
  one_mul   := Bool.true_and
  mul_one   := Bool.and_true
  mul_assoc := Bool.and_assoc
\end{lstlisting}

\paragraph{Haskell.}
\begin{lstlisting}[style=haskell]
-- (All, getAll) from Data.Monoid: mempty = True, mappend = (&&)
newtype ColdChain = ColdChain Bool

instance Monoid ColdChain where
  mempty = ColdChain True
  mappend (ColdChain a) (ColdChain b) = ColdChain (a && b)
\end{lstlisting}


% --- licListMonoid -------------------------------------------------
\subsection{Definition: \lean{licListMonoid}}

\begin{defn}[Licence-accumulation monoid]
For any licence type $L$, the \textbf{licence-list monoid} $(\mathrm{List}(L),
\mathbin{++}, [])$ accumulates the union (as a list) of licences required by
each leg along a route.
\end{defn}

\paragraph{Lean.}
\begin{lstlisting}[style=lean]
def licListMonoid (L : Type) : Monoid' where
  M         := List L
  mul       := (· ++ ·)
  one       := []
  one_mul   := List.nil_append
  mul_one   := List.append_nil
  mul_assoc := List.append_assoc
\end{lstlisting}

\paragraph{Haskell.}
\begin{lstlisting}[style=haskell]
-- Standard list monoid: mempty = [], mappend = (++)
type LicList l = [l]
\end{lstlisting}


%===================================================================
\section{Cold-chain worked example}
%===================================================================

\subsection{Example: warm last-mile leg breaks the cold chain}

\begin{prop}[Warm leg breaks the chain]
In a 2-node graph (\lean{Hub} $\to$ \lean{Shop}) with a single leg
\lean{lastMile} labelled $\mathrm{false}$ (warm), the tracking value of the
route Hub $\to$ Shop is $\mathrm{false}$, and the empty route is $\mathrm{true}$.
\end{prop}

\paragraph{Lean.}
\begin{lstlisting}[style=lean]
-- The cold-chain tracking induced by warmLabel (lastMile => false):
example : coldTracking.d Hub
    <|Shop, Path.cons (G := coldGraph) .lastMile
                      (.nil (G := coldGraph) Shop)|> = false

-- The empty route is intact:
example : coldTracking.d Hub <|Hub, .nil (G := coldGraph) Hub|> = true
\end{lstlisting}

\paragraph{Haskell.}
\begin{lstlisting}[style=haskell]
warmLabel :: Arc -> Bool
warmLabel LastMile = False

coldTracking :: Tracking ColdGraph Bool
coldTracking = ofEdge warmLabel

-- trackVal coldTracking Hub [LastMile] = False
-- trackVal coldTracking Hub []        = True
\end{lstlisting}


%===================================================================
\section{Translation conventions}
%===================================================================

\begin{enumerate}
\item \textbf{Bundled monoid.}  Lean uses a record \lean{Monoid'} instead of a
  typeclass, matching the bundled style of the rest of the formalisation.
  In Haskell, the standard \hask{Monoid} typeclass plays the same role.

\item \textbf{T-root and T-shift as data.}  The tracking laws are bundled into
  the \lean{Tracking} structure as proof fields, not stated as separate
  theorems.  In Haskell they are law comments; QuickCheck could test them.

\item \textbf{Proof bodies omitted.}  Only definitions and theorem
  \emph{statements} are translated.  Lean checks the proofs; the human
  verifier checks that the \emph{statements} match the intended mathematics.
\end{enumerate}

\end{document}
